% Intended LaTeX compiler: pdflatex
\documentclass[presentation]{beamer}
\usepackage[utf8]{inputenc}
\usepackage[T1]{fontenc}
\usepackage{graphicx}
\usepackage{longtable}
\usepackage{wrapfig}
\usepackage{rotating}
\usepackage[normalem]{ulem}
\usepackage{amsmath}
\usepackage{amssymb}
\usepackage{capt-of}
\usepackage{hyperref}
\usetheme{default}
\author{Lenin G. Falconi}
\date{2026-04-08}
\title{S1 Sistemas de Numeración}
\hypersetup{
 pdfauthor={Lenin G. Falconi},
 pdftitle={S1 Sistemas de Numeración},
 pdfkeywords={},
 pdfsubject={},
 pdfcreator={Emacs 29.3 (Org mode 9.6.15)}, 
 pdflang={Spanish}}
\begin{document}

\maketitle
\begin{frame}{Outline}
\setcounter{tocdepth}{1}
\tableofcontents
\end{frame}


\begin{frame}[label={sec:org96615ff}]{Sistemas de Numeración}
\begin{block}{Concepto General de Número}
\begin{itemize}
\item Un número se puede representar como una cadena de dígitos
\(b_{n-1} \dots b_2b_1b_0 \cdot b_{-1}b_{-2}\dots b_m\)
\item La posición relativa \(i\) del dígito posee un \alert{peso} \(r^i\), donde \(r\)
es la \alert{base} del sistema de numeración
\item El alfabeto de dígitos posibles en una base \(r\) cumple:
\(0 \leq b_i \leq r-1\)
\end{itemize}

Por tanto, un número \(N\) en base \(r\) se expresa como:

$$N_r = \sum_{i=-m}^{n-1} b_i \times r^i$$
\end{block}

\begin{block}{Tipos de Sistemas de Numeración}
\begin{center}
\begin{tabular}{lr}
Denominación & Base\\[0pt]
\hline
Decimal & 10\\[0pt]
Octal & 8\\[0pt]
Binario & 2\\[0pt]
Hexadecimal & 16\\[0pt]
\end{tabular}
\end{center}
\end{block}

\begin{block}{Ejercicio 1: Conversión a Decimal}
\begin{itemize}
\item Escriba un número en base 7
\item Escriba un número en base 4
\item ¿Qué valor decimal corresponde a cada uno?
\end{itemize}
\end{block}
\end{frame}

\begin{frame}[label={sec:org0353fa6}]{Sistema Decimal y Binario}
\begin{block}{Sistema de Numeración Decimal}
\begin{itemize}
\item La base es \(r=10\)
\item Los dígitos cumplen \(0 \leq b_i \leq 9\) y \(b_i \in \mathbb{Z}^+\)
\end{itemize}

Ejemplo:

$$123.45 = 1\times10^2 + 2\times10^1 + 3\times10^0 + 4\times10^{-1} + 5\times10^{-2}$$
\end{block}

\begin{block}{Sistema de Numeración Binario}
\begin{itemize}
\item La base es \(r=2\)
\item Los dígitos cumplen \(0 \leq b_i \leq 1\) y \(b_i \in \mathbb{Z}^+\)
\item Equivalencias directas: \(0_2 = 0_{10}\) y \(1_2 = 1_{10}\)
\end{itemize}

Ejemplo:

$$1101.01_2 = 1\times2^3 + 1\times2^2 + 0\times2^1 + 1\times2^0 + 0\times2^{-1} + 1\times2^{-2}$$
\end{block}
\end{frame}

\begin{frame}[label={sec:orgddc683a}]{Conversión de Base r a Decimal}
\begin{block}{Conversión \(N_r \rightarrow X_{10}\)}
Sea \(N_r \rightarrow X_{10}\) la conversión de un número en base \(r\) a decimal.

$$X_{10} = \sum_i b_i \times r^i$$
\end{block}

\begin{block}{Ejercicio 2: \(N_r \rightarrow X_{10}\)}
\begin{itemize}
\item \((4021.2)_5\)
\item \((127.4)_8\)
\item \((110101)_2\)
\item \((B65F)_{16}\)
\end{itemize}
\end{block}
\end{frame}

\begin{frame}[label={sec:org786ab80}]{Conversión Decimal a Binario}
\begin{block}{Parte Entera}
Dado un número decimal, se aplica divisiones sucesivas entre 2:

\begin{itemize}
\item Sea \(N\) la parte entera decimal
\item Dividir sucesivamente entre 2 y registrar residuos
\item El número binario buscado = residuos leídos en \alert{orden inverso}
\end{itemize}

$$N = 2N_1 + R_0 \quad N_1 = 2N_2 + R_1 \quad \dots \quad N_{m-1} = 2N_m + R_{m-1}$$

Entonces:

$$N = 2^mN_m + 2^{m-1}R_{m-1} + \dots + 2^1R_1 + R_0$$
\end{block}

\begin{block}{Ejercicio 3: \(N_{10} \rightarrow X_r\)}
\begin{enumerate}
\item Convertir 41 a binario
\item Convertir 153 a octal
\item Convertir 256 a hexadecimal
\end{enumerate}
\end{block}

\begin{block}{Parte Fraccionaria}
Sea \(F = 0\cdot b_{-1}b_{-2}\dots\) la parte fraccionaria en binario.

$$2F = b_{-1} + F_1$$

Procedimiento:
\begin{enumerate}
\item Multiplique la fracción por la base y tome la \alert{parte entera}
\item Repita con la nueva fracción resultante
\item Detenga al llegar a 0 o alcanzar la precisión deseada
\end{enumerate}
\end{block}

\begin{block}{Ejercicio 4: Parte Fraccionaria}
\begin{enumerate}
\item \((0.6875)_{10} \rightarrow X_2\)
\item \((0.513)_{10} \rightarrow X_8\)
\end{enumerate}
\end{block}
\end{frame}

\begin{frame}[label={sec:org7e038f4}]{Conversión entre Binario, Octal y Hexadecimal}
\begin{block}{Sistema Hexadecimal}
La base es 16 y sus símbolos son:

$$0,1,2,3,4,5,6,7,8,9,A,B,C,D,E,F$$

Relaciones clave:

$$2^3 = 8 \qquad 2^4 = 16$$
\end{block}

\begin{block}{Conversión Binario \(\rightarrow\) Octal}
\begin{enumerate}
\item Agrupar en bloques de \alert{3 bits} desde el punto decimal rellenando con 0
\end{enumerate}

Sea \(N_2 = 10101001.10111\):

$$010\,101\,001.101\,110$$

\begin{enumerate}
\item Sustituir cada grupo por su equivalente octal
\end{enumerate}

$$\Rightarrow 251.56_8$$
\end{block}

\begin{block}{Conversión Binario \(\rightarrow\) Hexadecimal}
\begin{enumerate}
\item Agrupar en bloques de \alert{4 bits} desde el punto decimal rellenando con 0
\end{enumerate}

Sea \(N_2 = 10101001.10111\):

$$1010\,1001.1011\,1000$$

\begin{enumerate}
\item Sustituir cada grupo por su equivalente hexadecimal
\end{enumerate}

$$\Rightarrow A9.B8_{16}$$
\end{block}
\end{frame}

\begin{frame}[label={sec:orge49928f},fragile]{Ordenes de Magnitud}
 \begin{block}{Potencias de 2 para Unidades de Información}
\begin{center}
\begin{tabular}{lll}
Potencia & Nombre & Valor\\[0pt]
\hline
\(2^0\) & bit & 1\\[0pt]
\(2^3\) & byte & 8\\[0pt]
\(2^{10}\) & Kilo & 1 024\\[0pt]
\(2^{20}\) & Mega & 1 048 576\\[0pt]
\(2^{30}\) & Giga & 1 073 741 824\\[0pt]
\(2^{40}\) & Tera & \textasciitilde{}1.1 billones\\[0pt]
\end{tabular}
\end{center}
\end{block}

\begin{block}{Ejercicio 5: Órdenes de Magnitud}
\begin{itemize}
\item Verifique con Python los valores de la tabla anterior
\end{itemize}

\begin{verbatim}
for nombre, pot in zip(['bit','byte','Kilo','Mega','Giga','Tera'],
                       [0, 3, 10, 20, 30, 40]):
    print(f"2^{pot:2d} -> {nombre:5s} = {2**pot}")
\end{verbatim}

\begin{verbatim}
2^ 0 -> bit   = 1
2^ 3 -> byte  = 8
2^10 -> Kilo  = 1024
2^20 -> Mega  = 1048576
2^30 -> Giga  = 1073741824
2^40 -> Tera  = 1099511627776
\end{verbatim}
\end{block}
\end{frame}

\begin{frame}[label={sec:org11413e2},fragile]{Taller}
 \begin{block}{Ejercicios Propuestos}
\begin{enumerate}
\item Función en Python: convierte un número en base \(r\) a decimal
\item Función en Python: convierte decimal a binario
\begin{itemize}
\item separar parte entera y fraccionaria
\end{itemize}
\item Generalizar para convertir de decimal a base \(r\) arbitraria
\end{enumerate}
\end{block}

\begin{block}{Plantilla Taller}
\begin{verbatim}
def numero_r_a_dec(numero: str, base: int):
    resp = 0
    # implementar
    print(f"{numero} en base {base} = {resp} en base 10")

def decimal_a_binario(numero: float):
    entero = int(numero)
    fraccion = numero - entero
    # implementar parte entera y fraccionaria
    pass

def decimal_a_base_r(numero: float, base: int):
    # generalizar decimal_a_binario para base r
    pass
\end{verbatim}
\end{block}
\end{frame}
\end{document}
